\scanpage{166}
Si
\[
d'_0=\operatorname{card}I'_0,
\]
\unsure{d} est la codimension (p.r. à $\deleted{X}\inserted{D_0}=\mathbb C^I$) du topos élémentaire envisagé, \deleted{qui est}\deleted{ici de $X$} alors que
\[
\unsure{d'_1}=\operatorname{card}I'_1
\]
est la dimension de ``l'âme'' du topos, qui est par définition
\[
\deleted{\unsure{D_{\pi}^*|D_{I'_1}^*}}=
\underbrace{\prod_{i\in I'_1}\mathbb C^*}_{\mathbb C^{\unsure{*}I'_1\unsure{*}}}
\times
\underbrace{\prod_{i\in I\setminus I'_1}\mathbb P}_{\mathbb P^{I'_0}}
\]
\margin{\deleted{$\pi$ désigne la partition envisagée de $I$}}
\deleted{On regarde $D_{I'_1}$ comme plongé dans}
\deleted{$\unsure{D_{I'_1}\times}$}
\deleted{On pose, pour tte partie $J$ de $I$}
\diagram{\begin{tikzcd}[ampersand replacement=\&,row sep=large]
\deleted{D_J=\mathbb C^{*J}\times\mathbb P^{I\setminus J}\subset X=\mathbb C^J}
\arrow[d,phantom,"\deleted{\cup}" description]\\
\deleted{D_J^*=\mathbb C^{\unsure{+|*}J}\times\unsure{\mathbb R|\mathbb P}^{I\setminus J}}
\end{tikzcd}}

Pour tte partie $J$ de $I$, on pose
\[
(269)
\]
\diagram{\begin{tikzcd}[ampersand replacement=\&,row sep=large]
D_J=\mathbb P^J\times\mathbb C^{I\setminus J}\subset\mathbb C^I=X
\arrow[d,phantom,"\cup" description]\\
D_J^*=\mathbb P^J\times\mathbb C^{*I\setminus J}
\end{tikzcd}}
\margin{sous-espace \unsure{typique|topique} défini par l'annulation des coordonnées $z_j$ \unsure{où|pour} $j\in J$, de codim $d=\operatorname{card}J$}

un ouvert dense
\[
(270)\qquad\left\{\begin{array}{l}
D_d=\displaystyle\coprod_{\substack{J\unsure{\subset|\in}P(I)\\\operatorname{card}J=d}}D_J\\
D_d^{\unsure{d|*}}=\displaystyle\coprod_JD_J^*
\end{array}\right.
\]
\margin{165}
